% Generado por generar_capitulos.py a partir de ../.tex8
\chapter*{Resumen}
\markboth{chapter name}{Sinopsis}

% <- resumen de main.tex
Un grupo que trabaja junto puede recordar más que cada uno de sus miembros por separado: cada
miembro sabe quién sabe qué y a quién preguntar.
Wegner \cite{wegner1987transactive} llamó a esta organización \emph{memoria
transactiva}.
Este trabajo pregunta si puede construirse solo con memorias asociativas, sin
clasificadores supervisados, y presenta EAM-TMS, un sistema distribuido de memoria
transactiva construido con las Memorias Asociativas Entrópicas de Pineda y Morales en su
variante ponderada (W-EAM) \cite{pineda2021entropic,pineda2022weighted} y su extensión
heteroasociativa (EHAM) \cite{morales2025missing}.
El sistema tiene ocho módulos especialistas, uno por cada categoría del conjunto de datos
ETH-80 \cite{leibe2003eth80}
(\texttt{apple}, \texttt{car}, \texttt{cow}, \texttt{cup}, \texttt{dog}, \texttt{horse},
\texttt{pear}, \texttt{tomato}).
Cada módulo guarda su contenido en una sola memoria heteroasociativa entre palabras e
imágenes y lleva dos directorios, uno textual y uno visual, con las transacciones en que
participa.
La arquitectura modela los tres mecanismos de Wegner \cite{wegner1995computer}: la
actualización del directorio, la asignación de información y la coordinación de la
recuperación.
Un mediador central coordina la fase temprana y deja de ser necesario en la fase madura,
cuando los directorios de los módulos bastan para enrutar.
Las memorias operan en los dos sentidos, de palabras a imágenes y de imágenes a
palabras.
Las palabras entran como vectores fastText de 300 dimensiones
\cite{bojanowski2017fasttext}, con vocabulario extraído de ConceptNet 5.7
\cite{speer2017conceptnet} y \emph{cuantización por magnitud} en 16 niveles; las
imágenes, como códigos de un ResNet18 \cite{he2016resnet}.
En el estudio de calibración, con un directorio compartido y un banco de 411 consultas
de texto ($N{=}400$ consultas de formación, 5 semillas), el 90.0\% de las consultas
llega al módulo correcto en la fase temprana y el 93.0\% en la fase madura con la
lectura del directorio normalizada por número de registros (B1); la cifra madura se
mide sobre las mismas consultas que formaron el directorio.
La lectura B1 reduce la participación del módulo \texttt{apple} del 20.5\% de la
lectura cruda al 13.8\% (ideal 12.5\%).
Con imágenes, los directorios visuales envían al especialista correcto el 96.2\% de las
656 imágenes de prueba (631), rechazan el 3.5\% y se equivocan en dos.
Con objetos que ningún módulo vio, el 89.2\% de las vistas llega al especialista y
recibe respuesta, frente al 95.5\% con vistas nuevas de objetos conocidos, con dos
enrutamientos falsos en 3280.
Frente a una sola memoria con las ocho clases, la partición en especialistas recuerda con
más fidelidad y, si además el módulo destino debe contener la consulta, acepta menos
consultas fuera de dominio cuando el contenido crece (2 frente a 8 de 40 con 200
imágenes por clase).
Los límites observados provienen de las representaciones (el léxico compartido entre
clases vecinas y la cobertura del espacio latente ante objetos nuevos), no del mecanismo
asociativo.

\hfill

\textbf{Palabras clave}: [PENDIENTE]
